\begin{textsample}{POS Dim 1 – 2009 – Score 23.00 – 2009/t000128.txt}  \label{ex:f1_pos_031}
So, my question : are we alone? \textbf{The} story of humans is \textbf{the} story of ideas — scientific ideas that shine light into dark corners, ideas that we embrace rationally and irrationally, ideas for which we ' ve lived and died and killed and been killed, ideas that have vanished in history, and ideas that have been set in dogma.
It ' s a story of nations, of ideologies, of territories, and of conflicts among them.
But, every moment of human history, from \textbf{the} Stone Age to \textbf{the} Information Age, from Sumer and Babylon to \textbf{the} iPod and celebrity gossip, they ' ve all been carried out — every book that you ' ve read, every poem, every laugh, every tear — they ' ve all happened here.
Here.
Here.
Here.
Perspective is a very powerful thing.
Perspectives can change.
Perspectives can be altered.
From my perspective, we live on a fragile \textbf{island} of life, in a \textbf{universe} of possibilities.
For many millennia, humans have been on a journey to find answers, answers to questions about naturalism and transcendence, about who we are and why we are, and of course, who else might be out there.
Is it really just us?
Are we alone in this vast \textbf{universe} of energy and matter and \textbf{chemistry} and physics?
Well, if we are, it ' s an awful waste of space.
But, what if we ' re not?
What if, out there, others are asking and answering similar questions?
What if they look up at \textbf{the} night sky, at \textbf{the} same stars, but from \textbf{the} opposite side?
Would \textbf{the} discovery of an older cultural civilization out there inspire us to find ways to survive our increasingly uncertain technological adolescence?
Might it be \textbf{the} discovery of a distant civilization and our common cosmic origins that finally drives home \textbf{the} message of \textbf{the} bond among all humans?
Whether we ' re born in San Francisco, or Sudan, or close to \textbf{the} heart of \textbf{the} Milky Way galaxy, we are \textbf{the} products of a billion-year lineage of wandering stardust.
We, all of us, are what happens when a primordial mixture of hydrogen and helium \textbf{evolves} for so long that it begins to ask where it came from.
Fifty years ago, \textbf{the} journey to find answers took a different path and SETI, \textbf{the} Search for Extra-Terrestrial Intelligence, began.
So, what exactly is SETI?
Well, SETI uses \textbf{the} tools of astronomy to try and find evidence of someone else ' s technology out there.
Our own technologies are visible over interstellar distances, and theirs might be as well.
It might be that some massive network of communications, or some shield against asteroidal impact, or some huge astro-engineering project that we can ' t even begin to conceive of, could generate signals at radio or optical \textbf{frequencies} that a determined program of \textbf{searching} might \textbf{detect}.
For millennia, we ' ve actually turned to \textbf{the} priests and \textbf{the} philosophers for guidance and instruction on this question of whether there ' s intelligent life out there.
Now, we can use \textbf{the} tools of \textbf{the} 21st century to try and observe what is, rather than ask what should be, believed.
SETI doesn ' t presume \textbf{the} existence of extra terrestrial intelligence ; it merely notes \textbf{the} possibility, if not \textbf{the} \textbf{probability} in this vast \textbf{universe}, which seems fairly uniform. \textbf{The} numbers suggest a \textbf{universe} of possibilities.
Our sun is one of 400 billion stars in our galaxy, and we know that many other stars have planetary systems.
We ' ve discovered over 350 in \textbf{the} last 14 years, including \textbf{the} small \textbf{planet}, announced earlier this week, which has a radius just twice \textbf{the} size of \textbf{the} Earth.
And, if even all of \textbf{the} planetary systems in our galaxy were devoid of life, there are still 100 billion other galaxies out there, altogether $10^{22}$ stars.
Now, I ' m going to try a trick, and recreate an experiment from this morning.
Remember, one billion?
But, this time not one billion dollars, one billion stars.
Alright, one billion stars.
Now, up there, 20 feet above \textbf{the} stage, that ' s 10 trillion.
Well, what about $10^{22}$?
Where ' s \textbf{the} line that marks that?
That line would have to be 3. 8 million miles above this stage. 16 times farther away than \textbf{the} moon, or four percent of \textbf{the} distance to \textbf{the} sun.
So, there are many possibilities.
And much of this vast \textbf{universe}, much more may be habitable than we once thought, as we study extremophiles on Earth — organisms that can live in conditions totally inhospitable for us, in \textbf{the} hot, high pressure thermal vents at \textbf{the} bottom of \textbf{the} ocean, frozen in ice, in boiling battery acid, in \textbf{the} cooling waters of nuclear reactors.
These extremophiles tell us that life may exist in many other environments.
But those environments are going to be widely spaced in this \textbf{universe}.
Even our nearest star, \textbf{the} Sun — its \textbf{emissions} suffer \textbf{the} tyranny of light speed.
It takes a full eight minutes for its radiation to reach us.
And \textbf{the} nearest star is 4. 2 light years away, which means its light takes 4. 2 years to get here.
And \textbf{the} edge of our galaxy is 75, 000 light years away, and \textbf{the} nearest galaxy to us, 2. 5 million light years.
That means any signal we \textbf{detect} would have started its journey a long time ago.
And a signal would give us a glimpse of their past, not their present.
Which is why Phil Morrison calls SETI,"\textbf{the} archaeology of \textbf{the} future."It tells us about their past, but detection of a signal tells us it ' s possible for us to have a long future.
I think this is what David Deutsch meant in 2005, when he ended his Oxford TEDTalk by saying he had two principles he ' d like to share for living, and he would like to carve them on stone tablets. \textbf{The} first is that problems are inevitable. \textbf{The} second is that problems are soluble.
So, ultimately what ' s going to determine \textbf{the} success or failure of SETI is \textbf{the} longevity of technologies, and \textbf{the} mean distance between technologies in \textbf{the} cosmos — distance over space and distance over time.
If technologies don ' t last and persist, we will not succeed.
And we ' re a very young technology in an old galaxy, and we don ' t yet know whether it ' s possible for technologies to persist.
So, up until now I ' ve been talking to you about really large numbers.
Let me talk about a relatively small number.
And that ' s \textbf{the} length of time that \textbf{the} Earth was lifeless.
Zircons that are mined in \textbf{the} Jack Hills of western Australia, zircons taken from \textbf{the} Jack Hills of western Australia tell us that within a few hundred million years of \textbf{the} origin of \textbf{the} \textbf{planet} there was abundant water and perhaps even life.
So, our \textbf{planet} has spent \textbf{the} vast majority of its 4. 56 billion year history developing life, not anticipating its emergence.
Life happened very quickly, and that bodes well for \textbf{the} potential of life elsewhere in \textbf{the} cosmos.
And \textbf{the} other thing that one should take away from this chart is \textbf{the} very narrow range of time over which humans can claim to be \textbf{the} dominant intelligence on \textbf{the} \textbf{planet}.
It ' s only \textbf{the} last few hundred thousand years modern humans have been pursuing technology and civilization.
So, one needs a very deep appreciation of \textbf{the} diversity and incredible scale of life on this \textbf{planet} as \textbf{the} first step in preparing to make contact with life elsewhere in \textbf{the} cosmos.
We are not \textbf{the} pinnacle of \textbf{evolution}.
We are not \textbf{the} determined product of billions of years of \textbf{evolutionary} plotting and planning.
We are one outcome of a continuing adaptational process.
We are residents of one small \textbf{planet} in a corner of \textbf{the} Milky Way galaxy.
And Homo sapiens are one small leaf on a very extensive Tree of Life, which is densely populated by organisms that have been honed for survival over millions of years.
We misuse language, and talk about \textbf{the}"ascent"of man.
We understand \textbf{the} scientific basis for \textbf{the} interrelatedness of life but our ego hasn ' t caught up yet.
So this"ascent"of man, pinnacle of \textbf{evolution}, has got to go.
It ' s a sense of privilege that \textbf{the} natural \textbf{universe} doesn ' t share.
Loren Eiseley has said,"One does not meet oneself until one catches \textbf{the} reflection from an eye other than human."One day that eye may be that of an intelligent alien, and \textbf{the} sooner we eschew our narrow view of \textbf{evolution} \textbf{the} sooner we can truly explore our ultimate origins and destinations.
We are a small part of \textbf{the} story of cosmic \textbf{evolution}, and we are going to be responsible for our continued participation in that story, and perhaps SETI will help as well.
Occasionally, throughout history, this concept of this very large cosmic perspective comes to \textbf{the} \textbf{surface}, and as a result we see transformative and profound discoveries.
So in 1543, Nicholas Copernicus published"\textbf{The} Revolutions of Heavenly \textbf{Spheres},"and by taking \textbf{the} Earth out of \textbf{the} center, and putting \textbf{the} sun in \textbf{the} center of \textbf{the} solar system, he opened our eyes to a much larger \textbf{universe}, of which we are just a small part.
And that Copernican revolution continues today to influence science and philosophy and technology and theology.
So, in 1959, Giuseppe Coccone and Philip Morrison published \textbf{the} first SETI article in a refereed journal, and brought SETI into \textbf{the} scientific mainstream.
And in 1960, Frank Drake conducted \textbf{the} first SETI observation looking at two stars, Tau Ceti and Epsilon Eridani, for about 150 hours.
Now Drake did not discover extraterrestrial intelligence, but he learned a very valuable lesson from a passing aircraft, and that ' s that terrestrial technology can interfere with \textbf{the} \textbf{search} for extraterrestrial technology.
We ' ve been \textbf{searching} ever since, but it ' s impossible to overstate \textbf{the} magnitude of \textbf{the} \textbf{search} that remains.
All of \textbf{the} concerted SETI efforts, over \textbf{the} last 40-some years, are equivalent to scooping a single glass of water from \textbf{the} oceans.
And no one would decide that \textbf{the} ocean was without \textbf{fish} on \textbf{the} basis of one glass of water. \textbf{The} 21st century now allows us to build bigger glasses — much bigger glasses.
In Northern California, we ' re beginning to take observations with \textbf{the} first 42 telescopes of \textbf{the} Allen Telescope \textbf{Array} — and I ' ve got to take a moment right now to publicly thank Paul Allen and Nathan Myhrvold and all \textbf{the} TeamSETI members in \textbf{the} TED community who have so generously supported this research. \textbf{The} ATA is \textbf{the} first telescope built from a large number of small dishes, and hooked together with \textbf{computers}.
It ' s making silicon as important as aluminum, and we ' ll grow it in \textbf{the} future by adding more antennas to reach 350 for more sensitivity and leveraging Moore ' s law for more processing capability.
Today, our signal detection algorithms can find very simple artifacts and noise.
If you look very hard here you can see \textbf{the} signal from \textbf{the} Voyager 1 spacecraft, \textbf{the} most distant human \textbf{object} in \textbf{the} \textbf{universe}, 106 times as far away from us as \textbf{the} sun is.
And over those long distances, its signal is very faint when it reaches us.
It may be hard for your eye to see it, but it ' s easily found with our efficient algorithms.
But this is a simple signal, and tomorrow we want to be able to find more complex signals.
This is a very good year. 2009 is \textbf{the} 400th anniversary of Galileo ' s first use of \textbf{the} telescope, Darwin ' s 200th birthday, \textbf{the} 150th anniversary of \textbf{the} publication of"On \textbf{the} Origin of Species,"\textbf{the} 50th anniversary of SETI as a science, \textbf{the} 25th anniversary of \textbf{the} incorporation of \textbf{the} SETI Institute as a non-profit, and of course, \textbf{the} 25th anniversary of TED.
And next month, \textbf{the} Kepler Spacecraft will launch and will begin to tell us just how frequent Earth-like \textbf{planets} are, \textbf{the} targets for SETI ' s \textbf{searches}.
In 2009, \textbf{the} U.
N. has declared it to be \textbf{the} International Year of Astronomy, a global festival to help us residents of Earth rediscover our cosmic origins and our place in \textbf{the} \textbf{universe}.
And in 2009, change has come to Washington, with a promise to put science in its rightful position.
So, what would change everything?
Well, this is \textbf{the} question \textbf{the} Edge foundation asked this year, and four of \textbf{the} respondents said,"SETI."Why?
Well, to quote :"\textbf{The} discovery of intelligent life beyond Earth would eradicate \textbf{the} loneliness and solipsism that has plagued our species since its inception.
And it wouldn ' t simply change everything, it would change everything all at once."So, if that ' s right, why did we only capture four out of those 151 minds?
I think it ' s a problem of completion and delivery, because \textbf{the} fine print said,"What game-changing ideas and scientific developments would you expect to live to see?"So, we have a fulfillment problem.
We need bigger glasses and more hands in \textbf{the} water, and then working together, maybe we can all live to see \textbf{the} detection of \textbf{the} first extraterrestrial signal.
That brings me to my wish.
I wish that you would empower Earthlings everywhere to become active participants in \textbf{the} ultimate \textbf{search} for cosmic company. \textbf{The} first step would be to tap into \textbf{the} global brain trust, to build an environment where raw data could be stored, and where it could be accessed and \textbf{manipulated}, where new algorithms could be developed and old algorithms made more efficient.
And this is a technically creative challenge, and it would change \textbf{the} perspective of people who worked on it.
And then, we ' d like to augment \textbf{the} automated \textbf{search} with human insight.
We ' d like to use \textbf{the} pattern recognition capability of \textbf{the} human eye to find faint, complex signals that our current algorithms miss.
And, of course, we ' d like to inspire and engage \textbf{the} next generation.
We ' d like to take \textbf{the} materials that we have built for education, and get them out to students everywhere, students that can ' t come and visit us at \textbf{the} ATA.
We ' d like to tell our story better, and engage young people, and thereby change their perspective.
I ' m sorry Seth Godin, but over \textbf{the} millennia, we ' ve seen where tribalism leads.
We ' ve seen what happens when we divide an already small \textbf{planet} into smaller \textbf{islands}.
And, ultimately, we actually all belong to only one tribe, to Earthlings.
And SETI is a mirror — a mirror that can show us ourselves from an extraordinary perspective, and can help to trivialize \textbf{the} differences among us.
If SETI does nothing but change \textbf{the} perspective of humans on this \textbf{planet}, then it will be one of \textbf{the} most profound \textbf{endeavors} in history.
So, in \textbf{the} opening days of 2009, a visionary president stood on \textbf{the} steps of \textbf{the} U.
S.
Capitol and said,"We cannot help but believe that \textbf{the} old hatreds shall someday pass, that \textbf{the} lines of tribe shall soon dissolve, that, as \textbf{the} world grows smaller, our common humanity shall reveal itself."So, I look forward to working with \textbf{the} TED community to hear about your ideas about how to fulfill this wish, and in collaborating with you, hasten \textbf{the} day that that visionary statement can become a reality.
Thank you.

% matched lemmas: array, chemistry, computer, detect, emission, endeavor, evolution, evolutionary, evolve, fish, frequency, island, manipulate, object, planet, probability, search, sphere, surface, the, universe
\end{textsample}
